\section{Bayesian Networks: Paths and D-Separation}\label{seq:d_separation}

This is a short overview. A complete treatment is given by Koller and
Friedman \cite{koller2009probabilistic}.

\subsection*{Bayesian networks}

A Bayesian network is a directed acyclic graph (DAG) whose nodes are
random variables, together with the factorisation
\[
    p(v_1, \dots, v_n) = \prod_{k=1}^n p(v_k \mid \mathrm{pa}(v_k)),
\]
where $\mathrm{pa}(v_k)$ are the parents of $v_k$. The graph encodes
conditional independences, which d-separation lets us read directly from
it.

\subsection*{Paths and the three elementary connections}

A \emph{path} between two nodes is a sequence of adjacent nodes,
\textbf{regardless of the direction of the edges}. For example, in
$A \leftarrow B \rightarrow C$ there is a path from $A$ to $C$ even
though no arrow goes from $A$ to $C$. Each intermediate node of a path
takes one of three roles, depending on the two edges that meet at it.

\begin{enumerate}
    \item \textbf{Chain} $A \rightarrow B \rightarrow C$. $A$ influences
    $C$ through $B$.
    \item \textbf{Fork} (common cause) $A \leftarrow B \rightarrow C$.
    $A$ and $C$ are correlated because they share the cause $B$.
    \item \textbf{Collider} (common effect) $A \rightarrow B \leftarrow
    C$. $A$ and $C$ are independent causes of $B$.
\end{enumerate}

\subsection*{Blocked paths}

Let $\mathcal{S}$ be the set of observed (conditioned-on) variables. A
path is \textbf{blocked} by $\mathcal{S}$ if it contains at least one
intermediate node $B$ such that
\begin{itemize}
    \item $B$ is a chain or fork node and $B \in \mathcal{S}$; or
    \item $B$ is a collider and neither $B$ nor any of its descendants is
    in $\mathcal{S}$.
\end{itemize}
Otherwise the path is \textbf{active}. Two sets of variables are
\textbf{d-separated} by $\mathcal{S}$ if every path between them is
blocked. If $\mathcal{X}$ and $\mathcal{Y}$ are d-separated by
$\mathcal{S}$, then $\mathcal{X} \indep \mathcal{Y} \mid \mathcal{S}$ in
every distribution that factorises according to the graph.

\subsection*{The two mechanisms, with examples}

\paragraph{Chains and forks: screening off.}
Observing the middle node of a chain or a fork \emph{blocks} the path.

\emph{Chain.} Let $A$ = \enquote{it rains}, $B$ = \enquote{the grass is
wet}, $C$ = \enquote{someone slips on the grass}:
$A \rightarrow B \rightarrow C$. Without information about the grass,
knowing that it rains makes slipping more likely. But once we know
whether the grass is wet, the rain tells us nothing more about slipping:
all its influence goes through the wetness. Formally,
$p(C \mid A, B) = p(C \mid B)$. The observed $B$ \emph{screens off} $C$
from $A$.

\emph{Fork.} Let $B$ = \enquote{the patient has the flu}, $A$ =
\enquote{fever}, $C$ = \enquote{cough}: $A \leftarrow B \rightarrow C$.
Fever and cough are correlated in the population, because both are
symptoms of the flu. Among patients known to have the flu, observing a
fever does not make a cough more likely: $A \indep C \mid B$.

\paragraph{Colliders: explaining away.}
A collider works the other way round: the path is blocked as long as the
collider is \emph{unobserved}, and observing it (or one of its
descendants) \emph{opens} the path.

Let $A$ = \enquote{it rained}, $C$ = \enquote{the sprinkler was on},
$B$ = \enquote{the grass is wet}: $A \rightarrow B \leftarrow C$. Rain
and sprinkler are independent: $A \indep C$. But once we observe that the
grass is wet, learning that it rained makes the sprinkler \emph{less}
likely, because the rain already explains the wet grass. Observing the
common effect makes its causes dependent. This phenomenon, and only this
one, is called \textbf{explaining away}.

\begin{inwords}
Information flows through chains and forks unless the middle node is
known. It flows through a collider only if the collider (or something it
causes) is known. Two variables are conditionally independent when every
route between them is cut by one of these rules.
\end{inwords}

\subsection*{Application to the DKF}

In the graph of Figure~\ref{fig:dkf_model}, the latent states form a
chain $\mathbf{z}_1 \rightarrow \mathbf{z}_2 \rightarrow \cdots$, each
$\mathbf{z}_t$ has a child $\mathbf{x}_t$, and each input $\mathbf{u}_t$ is
a parent of $\mathbf{z}_t$.
\begin{itemize}
    \item Given $\mathbf{z}_{t-1}$, every path from $\mathbf{x}_{t-1}$,
    $\mathbf{z}_{t-2}$ or $\mathbf{u}_{t-1}$ to $\mathbf{z}_t$ passes
    through $\mathbf{z}_{t-1}$ as a chain or fork node, and is blocked.
    This gives the factorisation of the generative model and the
    structure of the posterior (Sections~\ref{seq:generative_model} and
    \ref{sec:posterior}).
    \item The future observation $\mathbf{x}_{t+1}$ is \emph{not}
    d-separated from $\mathbf{z}_t$ by $\mathbf{z}_{t-1}$: the path
    $\mathbf{z}_t \rightarrow \mathbf{z}_{t+1} \rightarrow
    \mathbf{x}_{t+1}$ is a chain whose middle node is unobserved. This is
    why the inference network needs $\mathbf{x}_{t:T}$.
    \item In the AR DKF, the added edge $\mathbf{x}_{t-1} \rightarrow
    \mathbf{z}_t$ is a path with no intermediate node, which nothing can
    block. This is why the AR DKF posterior depends on past observations
    (Proposition~\ref{prop:ar_posterior}).
\end{itemize}

